\documentclass[10pt,a4paper]{amsart}
\usepackage{amsmath,amssymb,amsthm,graphicx}
\usepackage{verbatim}
\newenvironment{reference}{\comment}{\endcomment}
%\newenvironment{reference}{}{}
\newenvironment{slogan}{\comment}{\endcomment}
\newenvironment{history}{\comment}{\endcomment}

% For commutative diagrams we use Xy-pic
\usepackage[all]{xy}

% We use 2cell for 2-commutative diagrams.
\xyoption{2cell}
% Theorem environments.
%
\theoremstyle{plain}
\newtheorem{theorem}[subsection]{Theorem}
\newtheorem{proposition}[subsection]{Proposition}
\newtheorem{lemma}[subsection]{Lemma}

\theoremstyle{definition}
\newtheorem{definition}[subsection]{Definition}
\newtheorem{example}[subsection]{Example}
\newtheorem{exercise}[subsection]{Exercise}
\newtheorem{situation}[subsection]{Situation}

\theoremstyle{remark}
\newtheorem{remark}[subsection]{Remark}
\newtheorem{remarks}[subsection]{Remarks}

\numberwithin{equation}{subsection}

% Macros
%
\def\lim{\mathop{\mathrm{lim}}\nolimits}
\def\colim{\mathop{\mathrm{colim}}\nolimits}
\def\Spec{\mathop{\mathrm{Spec}}}
\def\Hom{\mathop{\mathrm{Hom}}\nolimits}
\def\Ext{\mathop{\mathrm{Ext}}\nolimits}
\def\SheafHom{\mathop{\mathcal{H}\!\mathit{om}}\nolimits}
\def\SheafExt{\mathop{\mathcal{E}\!\mathit{xt}}\nolimits}
\def\Sch{\mathit{Sch}}
\def\Mor{\mathop{\mathrm{Mor}}\nolimits}
\def\Ob{\mathop{\mathrm{Ob}}\nolimits}
\def\Sh{\mathop{\mathit{Sh}}\nolimits}
\def\NL{\mathop{N\!L}\nolimits}
\def\CH{\mathop{\mathrm{CH}}\nolimits}
\def\proetale{{pro\text{-}\acute{e}tale}}
\def\etale{{\acute{e}tale}}
\def\QCoh{\mathit{QCoh}}
\def\Ker{\mathop{\mathrm{Ker}}}
\def\Im{\mathop{\mathrm{Im}}}
\def\Coker{\mathop{\mathrm{Coker}}}
\def\Coim{\mathop{\mathrm{Coim}}}

% Boxtimes
%
\DeclareMathSymbol{\boxtimes}{\mathbin}{AMSa}{"02}

%
% Macros for moduli stacks/spaces
%
\def\QCohstack{\mathcal{QC}\!\mathit{oh}}
\def\Cohstack{\mathcal{C}\!\mathit{oh}}
\def\Spacesstack{\mathcal{S}\!\mathit{paces}}
\def\Quotfunctor{\mathrm{Quot}}
\def\Hilbfunctor{\mathrm{Hilb}}
\def\Curvesstack{\mathcal{C}\!\mathit{urves}}
\def\Polarizedstack{\mathcal{P}\!\mathit{olarized}}
\def\Complexesstack{\mathcal{C}\!\mathit{omplexes}}
% \Pic is the operator that assigns to X its picard group, usage \Pic(X)
% \Picardstack_{X/B} denotes the Picard stack of X over B
% \Picardfunctor_{X/B} denotes the Picard functor of X over B
\def\Pic{\mathop{\mathrm{Pic}}\nolimits}
\def\Picardstack{\mathcal{P}\!\mathit{ic}}
\def\Picardfunctor{\mathrm{Pic}}
\def\Deformationcategory{\mathcal{D}\!\mathit{ef}}
\usepackage{cleveref}
\crefname{originalreference}{original reference}{original references}
\Crefname{originalreference}{Original reference}{Original references}
\makeatletter
\newlabel{categories-lemma-extend-functor-by-colim}{{0FWY}{source}{}{sourcepointer-1}{}}
\newlabel{categories-lemma-extend-functor-by-colim@cref}{{[originalreference][0][]0FWY}{[1][0][]source}{}{}{}}
\makeatother
\crefname{originalreference}{original reference}{original references}
\Crefname{originalreference}{Original reference}{Original references}
\makeatletter
\newlabel{modules-lemma-coherent-finite-presentation}{{01BW}{source}{}{sourcepointer-2}{}}
\newlabel{modules-lemma-coherent-finite-presentation@cref}{{[originalreference][0][]01BW}{[1][0][]source}{}{}{}}
\makeatother
\crefname{originalreference}{original reference}{original references}
\Crefname{originalreference}{Original reference}{Original references}
\makeatletter
\newlabel{modules-lemma-finite-presentation-quasi-compact-colimit}{{01BS}{source}{}{sourcepointer-3}{}}
\newlabel{modules-lemma-finite-presentation-quasi-compact-colimit@cref}{{[originalreference][0][]01BS}{[1][0][]source}{}{}{}}
\makeatother
\crefname{originalreference}{original reference}{original references}
\Crefname{originalreference}{Original reference}{Original references}
\makeatletter
\newlabel{properties-lemma-quasi-coherent-colimit-finite-type}{{01PG}{source}{}{sourcepointer-4}{}}
\newlabel{properties-lemma-quasi-coherent-colimit-finite-type@cref}{{[originalreference][0][]01PG}{[1][0][]source}{}{}{}}
\makeatother
\crefname{originalreference}{original reference}{original references}
\Crefname{originalreference}{Original reference}{Original references}
\makeatletter
\newlabel{modules-lemma-limits-colimits}{{01AH}{source}{}{sourcepointer-5}{}}
\newlabel{modules-lemma-limits-colimits@cref}{{[originalreference][0][]01AH}{[1][0][]source}{}{}{}}
\makeatother
\title{Functors on coherent modules extend uniquely to filtered-colimit-commuting functors on QCoh, preserving additivity and exactness properties}
\author{The Stacks Project Authors}
\date{stacks-project:functors.tex:local-snapshot-20261012; GNU FDL 1.2 or later}
\begin{document}
\maketitle
\noindent\textit{Editorial note.} All delivered mathematical text is the author's own, copied verbatim from the pinned Stacks Project chapter file functors.tex; no statement or proof step was rewritten, reordered or removed.\par\medskip
\noindent\textit{Editorial note.} Span recomputation by label pairing (LANE\_BRIEF 10.2): in this chapter the author's proof environment follows the statement environment as a standalone block; the spec covers the complete environment pair (statement span, proof span) as recorded in the pool candidate.\par\medskip
\noindent\textit{Editorial note.} Occupation scan: combined claim\_key + statement\_anchor + generic-line scan over all registration files (active and withdrawn) found 0 hits for functors.tex (fu-occupied.json, 2026-10-14); consistent with the duty manager's six-check dispatch.\par\medskip
\ifdefined\hypertarget\hypertarget{sourcepointer-1}{}\fi
\noindent\textit{Source reference.} Original reference 0FWY: sources/stacks-chow/tags.txt, source lines 19863--19863. Cross-chapter standard prerequisite cited by the author inside the counted span; it is used as a pointer and its own proof is not part of this item.\par\medskip
\ifdefined\hypertarget\hypertarget{sourcepointer-2}{}\fi
\noindent\textit{Source reference.} Original reference 01BW: sources/stacks-chow/tags.txt, source lines 1741--1741. Cross-chapter standard prerequisite cited by the author inside the counted span; it is used as a pointer and its own proof is not part of this item.\par\medskip
\ifdefined\hypertarget\hypertarget{sourcepointer-3}{}\fi
\noindent\textit{Source reference.} Original reference 01BS: sources/stacks-chow/tags.txt, source lines 1737--1737. Cross-chapter standard prerequisite cited by the author inside the counted span; it is used as a pointer and its own proof is not part of this item.\par\medskip
\ifdefined\hypertarget\hypertarget{sourcepointer-4}{}\fi
\noindent\textit{Source reference.} Original reference 01PG: sources/stacks-chow/tags.txt, source lines 2229--2229. Cross-chapter standard prerequisite cited by the author inside the counted span; it is used as a pointer and its own proof is not part of this item.\par\medskip
\ifdefined\hypertarget\hypertarget{sourcepointer-5}{}\fi
\noindent\textit{Source reference.} Original reference 01AH: sources/stacks-chow/tags.txt, source lines 1690--1690. Cross-chapter standard prerequisite cited by the author inside the counted span; it is used as a pointer and its own proof is not part of this item.\par\medskip
\begin{lemma}
\label{lemma-functor-coherent}
Let $X$ and $Y$ be Noetherian schemes. Let
$F : \textit{Coh}(\mathcal{O}_X) \to \textit{Coh}(\mathcal{O}_Y)$
be a functor. Then $F$ extends uniquely to a functor
$\QCoh(\mathcal{O}_X) \to \QCoh(\mathcal{O}_Y)$
which commutes with filtered colimits.
If $F$ is additive, then its extension commutes with arbitrary direct sums.
If $F$ is exact, left exact, or right exact, so is its extension.
\end{lemma}

\begin{proof}
The existence and uniqueness of the extension is a general fact, see
Categories, Lemma \ref{categories-lemma-extend-functor-by-colim}.
To see that the lemma applies observe that coherent modules
are of finite presentation
(Modules, Lemma \ref{modules-lemma-coherent-finite-presentation}) and hence
categorically compact objects of $\textit{Mod}(\mathcal{O}_X)$ by
Modules, Lemma \ref{modules-lemma-finite-presentation-quasi-compact-colimit}.
Finally, every quasi-coherent module is a filtered colimit
of coherent ones for example by
Properties, Lemma \ref{properties-lemma-quasi-coherent-colimit-finite-type}.

\medskip\noindent
Assume $F$ is additive. If $\mathcal{F} = \bigoplus_{j \in J} \mathcal{H}_j$
with $\mathcal{H}_j$ quasi-coherent, then
$\mathcal{F} = \colim_{J' \subset J\text{ finite}}
\bigoplus_{j \in J'} \mathcal{H}_j$.
Denoting the extension of $F$ also by $F$ we obtain
\begin{align*}
F(\mathcal{F})
& =
\colim_{J' \subset J\text{ finite}}
F(\bigoplus\nolimits_{j \in J'} \mathcal{H}_j) \\
& =
\colim_{J' \subset J\text{ finite}}
\bigoplus\nolimits_{j \in J'} F(\mathcal{H}_j) \\
& =
\bigoplus\nolimits_{j \in J} F(\mathcal{H}_j)
\end{align*}
Thus $F$ commutes with arbitrary direct sums.

\medskip\noindent
Suppose $0 \to \mathcal{F} \to \mathcal{F}' \to \mathcal{F}'' \to 0$
is a short exact sequence of quasi-coherent $\mathcal{O}_X$-modules.
Then we write $\mathcal{F}' = \bigcup \mathcal{F}'_i$ as the
union of its coherent submodules, see
Properties, Lemma \ref{properties-lemma-quasi-coherent-colimit-finite-type}.
Denote $\mathcal{F}''_i \subset \mathcal{F}''$ the image of $\mathcal{F}'_i$
and denote $\mathcal{F}_i = \mathcal{F} \cap \mathcal{F}'_i =
\Ker(\mathcal{F}'_i \to \mathcal{F}''_i)$. Then it is clear that
$\mathcal{F} = \bigcup \mathcal{F}_i$ and
$\mathcal{F}'' = \bigcup \mathcal{F}''_i$
and that we have short exact sequences
$$
0 \to \mathcal{F}_i \to \mathcal{F}_i' \to \mathcal{F}_i'' \to 0
$$
Since the extension commutes with filtered colimits we have
$F(\mathcal{F}) = \colim_{i \in I} F(\mathcal{F}_i)$,
$F(\mathcal{F}') = \colim_{i \in I} F(\mathcal{F}'_i)$, and
$F(\mathcal{F}'') = \colim_{i \in I} F(\mathcal{F}''_i)$.
Since filtered colimits are exact
(Modules, Lemma \ref{modules-lemma-limits-colimits}) we
conclude that exactness properties of $F$ are inherited by
its extension.
\end{proof}

\end{document}
